% !TeX root = ../../packet-and-principal-models.tex
\section{The unit twist and the principal curve image}
\begin{proposition}[A global unit twist and its coarse involution]
\label{pp:prop:unitfamily}
On the actual family $(\mathcal B,\iota,\lambda_0)$ of
Theorem~\ref{thm:principalfamily}, the homomorphism
$\lambda_1=\lambda_0\circ\iota(\varepsilon_+)$ is another
principal polarization over the same base $C'$.
After forgetting auxiliary level, the two polarized families are
nonisomorphic over the identity of $C'$, and twisting twice returns
their polarized-family isomorphism class. On the full coarse
component of Proposition~\ref{pp:prop:explicit}, this operation is
the nontrivial involution
\[
 \tau_+(w)_i=\frac{w_i}{\sigma_i(\varepsilon_+)},\qquad
 F_{+,i}(z)=\frac{\sigma_i(a^2/\varepsilon_+)F_{t,i}(z)}
                         {3F_{t,i}(z)+1}.
\]
The specified B/C fibers receive their second polarization classes
from Proposition~\ref{pp:prop:actualpolarizations}.
\end{proposition}
\begin{proof}
Global RM self-adjointness makes $\lambda_1$ symmetric.
Its Riemann form multiplies each real eigenspace block of the
original positive form by $\sigma_i(\varepsilon_+)>0$.
Since $\varepsilon_+^{-1}\in\OL$, it is a principal polarization.
This is a global algebraic polarization homomorphism; no choice
of a global ample line bundle is required. The exact marking
comparison in Proposition~\ref{pp:prop:markeddistinct} identifies
$z+z^{-1}$ with the fixed global $t$, so the B/C twists are literal
instances of the preceding proposition.

In the standard module $P_{\rm std}$ put
$D=\operatorname{diag}(\varepsilon_+^{-1},1)$.
It preserves the lattice and satisfies
$E_{\varepsilon_+}(Dx,Dy)=E_1(x,y)$.
Evaluation at $w$ after $D$ is evaluation at $w/\varepsilon_+$.
This gives the displayed period formula with its stated direction.
Conjugating $\left(\begin{smallmatrix}r&\delta s\\v/\delta&h\end{smallmatrix}\right)$
by $D$ gives
$\left(\begin{smallmatrix}r&\delta s/\varepsilon_+\\
\varepsilon_+v/\delta&h\end{smallmatrix}\right)$,
so $D$ normalizes $\Gamma_{\rm std}$. Moreover
$D^2=\varepsilon_+^{-1}\operatorname{diag}(\varepsilon_+^{-1},\varepsilon_+)$;
the second factor belongs to $\Gamma_{\rm std}$ and the scalar
acts trivially on the domain. This proves the coarse involution.
Equivalently, twisting twice is pullback by $\iota(\varepsilon_+)$.
The distinct B/C polarization classes witness nontriviality of
both the coarse action and the two families over the identity base.
\end{proof}

\begin{proposition}[The unit involution preserves the principal curve]
\label{pp:prop:unitcurve}
The involution $\tau_+$ preserves the entire image $V_{\rm std}$ of
\eqref{pp:eq:curve}, and its restriction to that image is nontrivial.
More precisely, let $a_\triangle,b_\triangle,c_\triangle$ be the
original triangle generators, so that
$A=a_\triangle^2$, $B=b_\triangle$, $C=c_\triangle^2$, and put
$g=BACc_\triangle$. For the actual coarse moduli map
$f_0(z)=[WF_t(z)]$ one has
\[
 \tau_+f_0(z)=f_0(g^{-1}z).
\]
\end{proposition}
\begin{proof}
The full-triangle equivariance of the actual constructed map is essential.
McMullen's Section~4 identifies the quadratic trace-field extension and
its nontrivial automorphism, which sends the triangle generators to
$(-a_\triangle,b_\triangle,-c_\triangle)$.
Thus the two extensions of each real embedding of $L$ give identical
projective matrices and ordered target vertices. The vertex-respecting
triangle maps in \cite[Thm.~6.3, Eq.~(6.1)]{McMullen2026} therefore agree
for the two lifts. Keeping one from each pair yields the six components
of the actual $F$, with equivariance for the full triangle group.
Conjugating target coordinates by $h_t$ preserves this equivariance.
This argument would not follow from arbitrary $\Delta_0$-equivariance.

Set
\[
 s_0=-2-3t+4t^2+t^3-t^4,\quad
 h=s_0(I+C_t),\quad k=B_tA_tC_th.
\]
Exact field arithmetic gives $s_0^2=2-t$ and integral inverses for
$s_0$ and $\varepsilon_+=4-t^2$. Consequently
$\det h=\det k=\varepsilon_+$ and $h,k\in\GL_2^+(\OL)$.
If $\theta=2\cos(\pi/42)$, Cayley--Hamilton gives
$I+c_\triangle^2=\theta c_\triangle$.
Therefore $h$ is projectively the normalized $c_\triangle$, and
$F_t(gz)=kF_t(z)$, componentwise. The first-component base matrix
is $h_{t,1}^{-1}k_1h_{t,1}$ projectively; it is not $k_1$ in the
original base marking.

For $D=\operatorname{diag}(\varepsilon_+^{-1},1)$ the explicit witness
satisfies
\[
 T^{-1}k^{-\mathsf T}T\in\GL_2(\OL),\qquad
 U=\operatorname{diag}(\delta,1)^{-1}
       DWkW^{-1}\operatorname{diag}(\delta,1)\in\SL_2(\OL).
\]
The exact coefficient arrays, integral inverses and determinant identities
are checked by \path{certificates/curve_twist_image_v138.py} with
\path{certificates/curve_twist_image_data.json}.
Only this fixed word is needed; no search-completeness assertion enters.
In particular $\gamma=DWkW^{-1}$ belongs to the specified
$\Gamma_{\rm std}$. Hence
\[
 [DWF_t(z)]=[\gamma Wk^{-1}F_t(z)]
          =[WF_t(g^{-1}z)].
\]
Since $g^{-1}$ maps $\HH$ bijectively to itself, this proves equality
of the entire images. They are the actual algebraic images by
Theorem~\ref{thm:principalfamily} and
Proposition~\ref{pp:prop:unitfamily}.
The moved B/C points in Proposition~\ref{pp:prop:actualpolarizations}
show that the restricted involution is nontrivial.
\end{proof}

\begin{proposition}[Two principal classes over each specified RM variety]
\label{pp:prop:sixprincipal}
For each $X=A,B,C$, the underlying marked variety $(B_X,\iota_L)$ has
exactly two RM-compatible principal-polarization isomorphism classes,
represented by $\lambda_X$ and
$\lambda_X\circ\iota_L(\varepsilon_+)$.
Their origin-preserving RM-polarized automorphism groups are $\mu_6$.
The resulting six marked objects are pairwise distinct and all lie
on $V_{\rm std}$.
\end{proposition}
\begin{proof}
Identify rational homology with the rank-one module over the chosen
$M$. This remains valid for the induced A type. Since $M/L$ is
quadratic, every $L$-linear rational endomorphism has the unique form
$m_a+m_b\,\overline{(\cdot)}$, for $a,b\in M$.
After complexification the first summand preserves each embedding line,
while the second exchanges every line with its conjugate.
Each CM type contains exactly one line from every conjugate pair;
if $b\ne0$, all its embedding values are nonzero. Thus the second
summand cannot preserve the Hodge decomposition or cancel with the
first. It follows that
\[
 \operatorname{End}^0_L(B_X)=M,\qquad
 \operatorname{End}_L(B_X)=R,
\]
where the integral equality uses Proposition~\ref{prop:survivingorder}
for all three selected field actions. The first equality is a
centralizer statement, not an assertion that the full endomorphism
algebra at A is $M$.

If $\lambda$ is another RM-compatible principal polarization, then
$a=\lambda_X^{-1}\lambda$ is an integral automorphism commuting
with $L$. Symmetry makes $a$ fixed by the original Rosati, hence
$a\in E_L$, and positivity gives $a\in E_L^+$.
Conversely each such unit supplies a compatible principal polarization.
The $L$-linear automorphisms are precisely $R^\times$, whose norms
are $E_L^2$. Therefore the classes are
$E_L^+/E_L^2=\{1,\varepsilon_+\}$ by
Propositions~\ref{pp:prop:markedpolarizations}
and~\ref{pp:prop:actualpolarizations}.
The same norm computation gives the stated automorphism groups.
Proposition~\ref{pp:prop:unpolarizeddistinct} excludes identifications
between different $X$, while Proposition~\ref{pp:prop:unitcurve}
realizes each second class on the same image at $g^{-1}z_X$.
\end{proof}

\begin{proposition}[Unmarked separation at the induced A fiber]
\label{pp:prop:unmarkeda}
The principally polarized varieties
$(B_A,\lambda_A)$ and $(B_A,\lambda_A\circ\iota_L(\varepsilon_+))$
are nonisomorphic after forgetting all CM and RM markings.
\end{proposition}
\begin{proof}
Use the actual selected generator $Z_A=(A_t+tI)/(t^2-1)$.
Its lower-left and lower-right entries are
\[
 c=\frac{t^3-1}{t(t^2-1)},\qquad
 d_0=\frac{-t^5+5t^3+t^2-5t}{t^2-1}.
\]
In $M=\Ql(z)$, set $w=(z-d_0)/c$.
The row $(w,1)$ is an eigenvector of $Z_A^{\mathsf T}$, so
$(x,y)\mapsto xw+y$ identifies the actual rational homology with $M$.
On the common principal basis of Proposition~\ref{pp:prop:explicit},
its lattice and polarization parameter are
\[
 I_A=\sum_{j=0}^5\mathbf Z\,awt^j
       +\sum_{j=0}^5\mathbf Z\,b(1+3w)t^j,\qquad
 \xi_A=-\frac{c}{z-\bar z}.
\]
The trace form $\operatorname{Tr}_{M/\Ql}(\xi_Ax\bar y)$
reproduces the actual principal alternating matrix $E$ exactly.

In $\operatorname{End}_{\Ql}(M)=\bigoplus_\sigma M\sigma$,
the summands have disjoint permutation supports on the complex
embedding lines. A nonzero summand is a Hodge endomorphism precisely
when its automorphism stabilizes the actual type
$\{1,5,17,19,29,31\}$. Its stabilizer is $\{1,29\}$, and therefore
\[
 \operatorname{End}^0(B_A)=M+M\sigma_{29}.
\]
In particular the full algebra, of rational dimension $24$, is used;
restricting to multiplication or to one semilinear summand would
not prove the assertion.

Let $U_1,\ldots,U_{24}$ be the rational homology matrices for
$z^j$ and $z^j\sigma_{29}$, $0\le j<12$, in the displayed basis
of $I_A$. Form the $144\times24$ matrix $C$ by flattening them.
The full integral endomorphism coefficient lattice is
\[
 \mathcal L=\{v\in\Ql^{24}:Cv\in\mathbf Z^{144}\}.
\]
Equivalently, it is the rational dual of the integer lattice
generated by the constraint rows. This gives a complete lattice
calculation, including cancellations between the two summands.

Let $Q_+$ be the homology action of $\varepsilon_+$.
An unmarked polarized isomorphism $U$ would satisfy
$U^{\mathsf T}EU=EQ_+$, hence
\[
 \operatorname{Tr}_{\Ql}(Q_+^{-1}U^\dagger U)=12,\qquad
 U^\dagger=E^{-1}U^{\mathsf T}E.
\]
This is a positive definite quadratic form on the real Hodge
endomorphism algebra: the positive Hodge metric makes $\dagger$
the adjoint, and $Q_+$ is positive self-adjoint.
Its exact Gram matrix on $\mathcal L$ has an integral unimodular
reduction $H$ with positive rational LDL data $(\mu,d)$.
Every vector in its closed norm-$12$ ball satisfies
$|x_i|^2\le12(H^{-1})_{ii}$. The exact decomposition
\[
 x^{\mathsf T}Hx=\sum_i d_i
       \left(x_i+\sum_{j>i}\mu_{ji}x_j\right)^2
\]
therefore supplies finite coordinate bounds and exhaustive
suffix pruning using only rational comparisons.

The certificate \path{certificates/a_unmarked_polarization_v140.py}
reconstructs the actual model, full dual lattice and Gram matrix,
checks the integral unimodular reduction, positive LDL data and
bounds in \path{certificates/a_unmarked_polarization_data.json},
and exhausts this ball in $347$ nodes: only zero occurs.
The separate implementation
\path{independent/a_unmarked_polarization_v140.py} works directly
on $L^2$ homology, constructs $\sigma_{29}$ by a cyclic vector,
proves fullness by integer row saturation, and recomputes the
Gram and LDL in reverse coordinate order. Its $925$ nodes again
leave only zero. Thus the required norm-$12$ vector does not exist.
\end{proof}

\begin{proposition}[The four B/C polarization pairs]
\label{pp:prop:polarizationsquare}
Put $B_e=(B_B,\lambda_B\varepsilon_+^e)$ and
$C_f=(B_C,\lambda_C\varepsilon_+^f)$, for $e,f\in\{0,1\}$.
For every pair $B_e\to C_f$ the least positive integer scalar
polarization multiplier is $2$, attained by exactly six isogenies,
all of degree $64$.
With $\gamma_0=z^6(1-z^2)$ and
$v_{ef}=\sigma_{31}(\gamma_0)^e\gamma_0^{-f}$,
their normalized Hom parameters are
\[
 \eta=\begin{cases}
 \alpha v_{ef}z^{7j},&e=f,\\
 \bar\alpha v_{ef}z^{7j},&e\ne f,
 \end{cases}\qquad 0\le j<6.
\]
The six maps share one kernel for each pair. The pairs $(0,0),(1,1)$
have the same kernel $K_+$, while $(0,1),(1,0)$ have a second kernel
$K_-$. These are complementary order-$64$ subgroups of $B_B[2]$.
\end{proposition}
\begin{proof}
Proposition~\ref{pp:prop:homideal} gives every rational homomorphism
as $F_\eta=(m_0/\alpha)\eta\sigma_{31}$, integral exactly when
$\eta\in\mathfrak a$. Its source pullback multiplier is
$\sigma_{19}(\eta\bar\eta)$. For the two twisted forms the required
multiplier-$n$ equation is consequently
\[
 \eta\bar\eta=n\frac{\sigma_{31}(\varepsilon_+)^e}
                         {\varepsilon_+^f}.
\]
Since $\gamma_0\bar\gamma_0=\varepsilon_+$,
$N_{M/L}(v_{ef})$ is the unit factor on the right.

The polynomial $\Phi_{42}$ modulo $2$ is the product of the two
irreducible sextics
\[
 (X^6+X^4+X^2+X+1)(X^6+X^5+X^4+X^2+1).
\]
Together with $\alpha\bar\alpha=2$ and
$\alpha+\bar\alpha=1$, this identifies the two conjugate primes
above $2$ as $(\alpha),(\bar\alpha)$.
Every integral solution for $n=2$ therefore has one of these two
ideals. Division by $\alpha v_{ef}$ or $\bar\alpha v_{ef}$ gives
an integral unit of relative norm one, hence a root of unity.
Thus the complete candidates are the two families
$\alpha v_{ef}z^j$, $\bar\alpha v_{ef}z^j$, $0\le j<42$.

For membership in $\mathfrak a=\alpha\OL+\mathfrak P$,
consider the quotient $\mathbf F_{729}^\times/\mathbf F_{27}^\times$
of order $28$. The class of $z$ has order $7$.
The class of $\gamma_0$ is the unique element of order two:
$\gamma_0/\bar\gamma_0=-z^{14}$ reduces to $-1$.
The class of $\bar\alpha/\alpha$ is that same element, since
$\alpha^2-\alpha+2=0$ modulo $3$ gives
$\alpha^4=-1$ and $(\alpha/\bar\alpha)^2=-1$ there.
The automorphism $\sigma_{31}$ fixes this unique order-two class.
Consequently the displayed branch cancels the order-two class
exactly as stated, and the remaining exponent must be divisible by $7$.
There are exactly six maps. Their full norm, ideal membership,
integral homology matrices and polarization identities are verified in
\path{certificates/bc_polarization_square_v141.py} with its retained data.

Multiplier one would imply an isomorphism of the underlying B/C
varieties, excluded by Proposition~\ref{pp:prop:unpolarizeddistinct}.
Thus the displayed multiplier two is minimal. Multiplication by
$z^{7j}\in R^\times$ preserves each target form, so the six maps
for a fixed pair have the same kernel.

Let $T_{ef}$ denote the integral homology matrix for $j=0$.
The certificate verifies $\det T_{ef}=64$ and the four actual
polarization equations. The polarized adjoint relation implies
that every kernel lies in $B_B[2]$. It is therefore identified with
$\ker(T_{ef}\bmod2)$ on the common source basis.
All four ranks are $6$; the stacked ranks for $(00,11)$ and
$(01,10)$ are $6$, and the stacked rank for $(00,01)$ is $12$.
These exact finite calculations prove the asserted equality and
complementarity of kernels.
\end{proof}

\begin{remark}[Image and level boundaries]\label{pp:rem:unitimage}
The two parametrized moduli maps differ at the specified B/C parameters,
but their images agree by Proposition~\ref{pp:prop:unitcurve}.
The diagonal matrix in the square calculation need not lie in a chosen
neat subgroup, so no involution preserving an arbitrary fine-level
marking is asserted. These constructions neither add a point to the
original trace-polarized packet nor identify a principal image with
Mostaed's literal normalization.
\end{remark}

\begin{corollary}[Six Weil points on the specified principal image]
\label{pp:cor:weilimage}
Let $\mathcal W_{K,\rm std}$ denote the intrinsic $K$-Weil locus in the
specified principal RM component. Then
$V_{\rm std}\cap\mathcal W_{K,\rm std}$ contains the six distinct
marked points of Proposition~\ref{pp:prop:sixprincipal}.
All have $K$-signature $(3,3)$ and Hermitian discriminant $-1$.
The actual CM types and Hodge algebras are those of their respective
A/B/C underlying varieties. All six remain distinct as principally polarized varieties after
forgetting the RM and CM markings.
\end{corollary}
\begin{proof}
Keep the original lattice $P$, chosen field actions and integral
$\mathcal O_K$ generators of Proposition~\ref{prop:commonprincipal}.
Twisting replaces $E_1$ by
$E_{\varepsilon_+}(v,w)=E_1(\varepsilon_+v,w)$, where
$\varepsilon_+I_2$ is an integral automorphism of $P$.
The chosen $M$ action and therefore the actual CM type are unchanged.
When the new point is represented at $g^{-1}z_X$, these are the
transported actions through the explicit orbit identification above;
they are not inferred from a companion CM type.

The integral $K$ action commutes with $\varepsilon_+$ and its Rosati
remains conjugation. The complex $K$ action is unchanged, so the
signature is still $(3,3)$. On the six-dimensional $K$ homology space,
$M=LK$ and $L\cap K=\Ql$ give
\[
 \det_K(m_{\varepsilon_+})
 =N_{M/K}(\varepsilon_+)=N_{L/\Ql}(\varepsilon_+)=1.
\]
The last norm is also checked in the exact transport certificate.
Thus twisting the Hermitian form in one variable leaves its determinant
and discriminant unchanged. Positivity and principality were established
in Proposition~\ref{pp:prop:unitfamily}.
Changing polarization does not change the Hodge structure, so the
previous divisor/Weil generator and algebraicity conclusions apply.
Finally, the four B/C polarized objects remain distinct without markings
by Proposition~\ref{pp:prop:actualpolarizations}, and
Proposition~\ref{pp:prop:unmarkeda} separates the two A objects.
The underlying A variety is different from both B/C varieties by
Proposition~\ref{pp:prop:unpolarizeddistinct}. Hence all six remain
distinct without markings; no classification of all A polarizations
is needed.
\end{proof}

\begin{remark}[The integral comparison still required]
\label{pp:rem:obstruction}
For any full $O$ lattice $P'\subset L^2$ self-dual for $E_1$, its
determinant fractional ideal is $\bigwedge_O^2P'=\mathfrak D_L^{-1}$.
Indeed at every finite prime a local basis matrix $T_{\mathfrak p}$
turns the pairing into the standard alternating form multiplied by
$\det T_{\mathfrak p}$; perfection of the trace pairing is equivalent
to $(\det T_{\mathfrak p})O_{\mathfrak p}
=\mathfrak D_{L,\mathfrak p}^{-1}$. These equalities determine the
fractional ideal globally. Proposition~\ref{prop:noprincipal} shows
that the codifferent has no totally positive generator. Thus no
$T'\in\GL_2^+(L)$ satisfies $T'O^2=P'$. The positive matrix $W$ above
compares $P$ to the different module $P_{\rm std}$ and does not contradict
this obstruction.

The original trace-polarized fibers do not become principal by passing
to a cover alone. The chosen quotient kernel has orbit size $28$, so its
marked descent requires its stabilizer; this does not rule out descent
of an unmarked target or a suitable unspecified finite-index subgroup.
In particular Corollary~\ref{pp:cor:weilimage} concerns the explicit
principal image just constructed. Identifying it with Mostaed's named
$V$ requires an integral polarization/level comparison and equality of
curve images or direct point membership. Neither that identification
nor disjointness of the two images follows here. The original
standard lattice, a principal isogenous model, and literal membership
in a prescribed principally polarized locus remain distinct assertions.
\end{remark}

